\newpage
\chapter{The Natural Numbers}

\subsection{Definition as Sets}

We define the members of \(\N\) as:

\begin{itemize}
	\item \(0 := \emptyset\)
	\item \(1:= \{0\} = \{\emptyset, \{\emptyset\}\}\)
	\item \(2:= \{0, 1\} = \{\emptyset, \{\emptyset, \{\emptyset\}\}\}\)
	\item \(n:= \{n_1, \dots, 0\}\)
\end{itemize}

\subsection{Addition}

\[
	S(n) := n \cup {n} := n + 1
\]

\subsection{Inductive Sets}

A set \(I\) is inductive if

\[
	0 \in I, x \in I \implies S(x) \in I.
\]

\subsection{Axiom of Infinity}

An inductive set exists.

\subsection{Definition of the Sets of Naturals Numbers}

\[
	\N := \{x \mid x \in I, \text{ for every inductive set }\}
\]

\textbf{Proof that \(\N\) is inductive for every inductive set \(I, \N \subseteq I\):}

\(\forall I, \emptyset \in I\), hence \(0 \in \N\), for \(n \in \N\), let \(I\) be inductive then 

\[
	S(n) \in I \therefore S(n) \in \N
\]

\subsection{The Standard Order}

We define \(<\) for the natural numbers as: 

\[
	\forall m, n \in \N, m < n \iff m \in n.
\]

\subsubsection*{Lemma}

\begin{itemize}
	\item \(0 < n, \forall n \in \N\)
	\item \(\forall k,n \in \N, k < S(n) \iff k < n \lor k = n\)
\end{itemize}

\textbf{Proof:}

Suppose for the sake of contradiction that \(\exists n \in \N, n \ne 0: n < 0\). By definition 
\(n < 0 \iff n \in 0 = \emptyset\). Hence, we arrive at our desired contradiction.

For the second part of the lemma \(k < S(n) \iff k < n \lor k = n\):

\begin{align*}
	k < S(n) \implies k \in \{n\} \cup n &\iff k \in n \lor k = n \\ 
										 &\iff k < n \lor k = n \\
										 &\iff k < n
\end{align*}

\QED

\subsection{The Peano Axioms}

We will define the set of natural numbers, \( \N \), via the following 9 axioms. 
These axioms are known as the \emph{Peano Axioms}. The first 4 axioms define equality on 
the set \( \N \).

\begin{enumerate}[label=\Roman*.]

	\item \emph{Axiom of Reflexivity}: For every \( x \in \N \), we have \( x = x \).
	
	\item \emph{Axiom of Symmetry}: For every \( x, y \in \N \), if \( x = y \) then \( y = x \).
	
	\item \emph{Axiom of Transitivity}: For every \( x, y, z \in \N \), if \( x = y \) and \( y = z \) then \( x = z \). 
	
	\item \emph{Axiom of the Closure of Equality}: For all \( x, y \), if \( x \in \N \) and \( x = y \), then \( y \in \N \).\hfill 

\end{enumerate}

The remaining 5 axioms define the structure of \( \N \):

\begin{enumerate}[label=\Roman*.]
	\setcounter{enumi}{4}

	\item \( 1 \in \N \)
	
	\item If \( x \in \N \), then the successor \( S(x) \in \N \).
	
	\item There is no \( x \in \N \) such that \( S(x) = 1 \).
	
	\item For all \( x, y \in \N \), if \( S(x) = S(y) \), then \( x = y \).
	
	\item Let \( P(x) \) be a statement about the natural number \(x\). If:
	
		\begin{itemize}
	
			\item \( P(1) \) is true, and
	
			\item for all \( n \in \N \), if \( P(n) \) is true, then \( P(S(n)) \) is also true,
	
		\end{itemize}
	
		then \( P(x) \) is true for all \( x \in \N \).\hfill (Mathematical Induction)
\end{enumerate}

As shorthand, we denote:

\[
	S(1) = 2, \quad S(S(1)) = 3, \quad S(S(S(1))) = 4, \quad \text{and so on.}
\]

\subsection{Propositions and Proofs}

\subsubsection{Proposition 1: \texorpdfstring{\(n \ne m \implies S(n) \ne S(m)\)}{n!= m implies S (n)!=S (m)}}

\textbf{Proof:} 

We will prove this by contradiction. Assume \( n \ne m \) and \( S(n) = S(m) \). 
By Axiom 8, we have \( n = m \), which is a contradiction. Therefore, \( S(n) \ne S(m) \).

\subsubsection{Proposition 2: \texorpdfstring{\text{For any} \(n \in \N,\ n \ne S(n)\)}{ For any n in N, n!= S (n)}}

\textbf{Proof:} 

\(M = \{n \in \N \| n \ne S(n) \}\ \)

By \(1 \ne S(n)\ \) for any \(n \in \N\), this implies that 1 is part of the set \(M\). 
This implies that \(S(n) \ne S(S(n)) \implies S(n) \in M\ \) By Axiom 9 \(M = \N\)

\subsubsection{Proposition 3: \texorpdfstring{\(n \ne 1\ \exists m \in \N \mid n = S(m)\)}{n!= 1, exists m in N | n = S (m)}}

\textbf{Proof:}

\[
	M = \{1\} \cup\{n \in \N | \text{Proposition 3 is true}\}
\]

We know that 1 ins in the set \(M\). And by proposition 1 we know that

\[
	S(n) = S(S(m)) \implies S(n) \in M
\]

And by Axiom 9 we know that \(M = \N\)

\subsection{Definition of Addition in \texorpdfstring{\(\N\)}{}}

For any pair n, m \(\in \N\) there is a unique way to define

\[
	\text{Add}(n , m) = n + m
\]

\begin{enumerate}
	
	\item \textbf{Base Case:} \(n + 1 = S(n)\)

	\item \textbf{Inductive Step:} \(n + S(m) = S(n + m) \iff S(n + m)\)

\end{enumerate}

\textbf{Uniqueness:} 

Suppose: \textit{A} \& \textit{B} satisfy our conditions. Fix n and then let \(M = \{ m \in \N | 
A(n,m) = B(n, m)\}\). Then

\[
	A(n,1) = S(n) = B(n,1) \implies 1 \in M
\]

\[
	m \in M \implies A(n, m) = B(n, m) \implies A(n, S(m)) = S(A(n, m)) = S(B(n, m)) = B(n, S(m))
\]

\[
	\implies A(n , S(m)) = B(n, S(m))
\]

by Axiom 9 we know that \(M = \N\) and A = B

\textbf{Construction:} For \(n = 1\) Define \(A(n, m) = S(m)\)

\begin{enumerate}
	
	\item \(A(n, 1) = S(n) = S(1)\)
	
	\item \(A(n, S(m)) = S(A(n, m)) = S(S(m))\)

\end{enumerate}

Define: \(A(S(n), S(m)) = S(A(n, m))\)

\begin{enumerate}

	\item \(A(S(n), 1) = S(A(n, 1)) = S(S(n))\)

	\item \(A(S(n), S(m)) = S(A(n, m)) = S(S(m))\)

\end{enumerate}

\textbf{Commutativity of Addition:}

The proposition says \(n + m = m + n\) for any \(n, m \in \N\). We will prove this by 
induction on \(m\). Fix n and consider \(M = \{ n \in \N | A(n, m) = A(n, m)\}\)

Now recall that \(A(n, 1) = S(n)\)

For \(n = 1\):

\[
	A(n, 1) = S(1) = A(1, n) \implies 1 \in \N
\]

also \(A(n , k) = 1 + k \implies 1 + m = S(m) \implies 1 + m = m + 1 \implies 1 \in \N\)

Suppose: \(n \in \N \implies n + m = m + n\) or \(A(n, m) = A(m, n)\)

By construction \(A(S(n), m) = S(A(n, m))\) and by definition
\(A(S(n), m) = S(A(n, m)) = A(m, S(n)) = S(n) + m =  m + S(n) \implies S(n) \in M \implies 
\text{by induction } M = \N\).
